%!TEX root = ../LLMSeminar.tex
% ──────────────────────────────────────────────────────────────
%  Section 13 — Seminar II discussion (delivered): Verifiability
%               and the limits of self-improvement
% ──────────────────────────────────────────────────────────────
%
%  This section is a faithful summary of the live discussion in
%  the second seminar.  It records material that arose in the
%  room and that complements, rather than duplicates, the
%  pre-prepared notes of \S\ref{sec:recap-agency}--\S\ref{sec:ambitious-science}.
%  It is intended to be folded back into those sections at the
%  author's discretion, or to remain as a standalone chapter on
%  the verification hierarchy and the singularity objection.
% ──────────────────────────────────────────────────────────────

\section{Verifiability and the limits of self-improvement}%
\label{sec:verification}

The intent of the second seminar was to extend the picture of
\S\ref{sec:agent-loop}--\S\ref{sec:self-improvement} in two
directions: outward, into the operational realities of running
these systems on the hardware participants actually have; and
upward, into the question of what it means to \emph{trust} the
output of a stochastic, high-throughput source.  The
discussion that resulted was substantially less linear than
the prepared notes, and the present section records what
emerged.

\subsection{The pre-bootstrap phase}\label{sec:prebootstrap}

We opened with a debrief of the second homework of
\S\ref{sec:bitter-lesson}: install Ollama and run a local
model.  Several participants had attempted this with very
mixed success.

\begin{audience}[Student reports]
  \begin{itemize}
    \item Gemma~4 in its $26\,\text{B}$-parameter variant
      (with $\sim4\,\text{B}$ \emph{active} parameters under
      the mixture-of-experts routing) was made to run on a
      laptop with $16\,\text{GB}$ of integrated graphics
      memory, after manually tuning the number of layers
      offloaded to the GPU.  Useful throughput was reached at
      $13$--$16$ tokens/s, and the model produced a usable
      \LaTeX\ rendering of an exercise sheet in roughly five
      minutes.
    \item Qwen~3.5 in the same regime, by contrast, fell into
      a thinking loop on most prompts---``overthinking,'' as
      the room agreed---and produced no usable output.  This
      pathology, of small reasoning models recursing into
      their own scratchpad, was reported as a recurring
      failure mode of the current generation.
    \item The smallest variant of Gemma~4 ($2\,\text{B}$
      effective parameters) cannot reproduce a passage from
      Harry Potter on demand, despite the larger versions
      doing so easily.  The room concluded, only half in
      jest, that the small models do not contain
      ``culture''---they have been distilled past the point
      at which idiosyncratic training data survives.
    \item One participant, calling a local model directly via
      the Hugging Face \texttt{transformers} library (i.e.\
      bypassing the Ollama harness), found that the model
      \emph{denied} that it was running locally and
      maintained that it was in the cloud, until the same
      claim was put in the role of a system prompt
      (an ``admin'' role), at which point it accepted the
      premise and proceeded.  This is a small but instructive
      data point on prompt hierarchies.
  \end{itemize}
\end{audience}

\begin{keyresult}[The pre-bootstrap phase]\label{kr:prebootstrap}
  Most participants are at present in a \emph{pre-bootstrap}
  regime: the local-model setup must be done manually because
  the locally-runnable models are not yet capable enough to
  set themselves up.  This is a transient state.  The interesting
  threshold is the moment at which a local model is good
  enough to act as the agent that installs and configures the
  next, better, local model on one's behalf.  Until that
  threshold is crossed, the cost of running locally is
  measured in afternoons of human typing.
\end{keyresult}

A practical consequence: the cost-effective workflow at
present is heterogeneous.  Specialised models do specialised
jobs (\texttt{whisper.cpp} for speech, Marker for
PDF-to-Markdown with embedded \LaTeX, locally-run small models
for short generative tasks); the LLM is invoked only for the
genuinely natural-language part of the pipeline.  Routing the
work to the cheapest competent model, rather than putting
everything through the largest available LLM, is the
operationally correct discipline.

\begin{remark}[On hardware]
  The hardware target most authors of these models optimise
  for is presently the $32\,\text{GB}$ Apple-silicon laptop:
  unified memory removes the CPU/GPU partitioning that makes
  laptop inference painful elsewhere.  An intermediate option
  for our group will be a desktop with an RTX~3090 (sufficient
  for Marker, Whisper, and small Qwen-class models), with a
  shared server in the basement targeted as the eventual
  ``decent'' tier.
\end{remark}

\subsection{The verification hierarchy}\label{sec:verification-hierarchy}

The thread connecting almost every difficulty reported in
\S\ref{sec:prebootstrap} is the same one that troubled us at
the end of Seminar~I: the output of an LLM is a
\emph{high-frequency, low-trust} channel.  We have plenty of
text; we cannot at the moment quite believe any of it.  How
much error can be tolerated, and what mechanisms reduce it,
is therefore the central operational question.

\begin{intuition}[Information-theoretic framing]
  An LLM is a noisy channel: it carries genuine signal at
  some rate together with errors at some rate.  Shannon
  teaches that a noisy channel can be made arbitrarily
  reliable by coding---most simply, by repetition.  If the
  channel produces the right answer with probability~$p > 1/2$
  on a binary decision, sampling $N$ times and taking a
  majority vote drives the error rate to zero exponentially
  in~$N$.  The fact that compute is cheap and inference is
  embarrassingly parallel makes this a serious option, not
  an asymptotic curiosity.
\end{intuition}

This motivates a hierarchy of verification mechanisms,
ordered by the strength of the guarantee each provides.

\begin{keyresult}[The verification hierarchy]\label{kr:verification}
  In ascending order of confidence:
  \begin{enumerate}[(1)]
    \item \textbf{Consensus.}  Ask several independent
      sources---models, journals, colleagues---and take the
      majority view.  Cheap, broad, and noisy.
    \item \textbf{Computational test.}  Have the candidate
      result emit a programme that runs and produces an
      answer one can compare against known cases.  Higher
      bar; it forces the result to pass through a concrete
      executable bottleneck.
    \item \textbf{Formalisation.}  Encode the statement and
      its proof in a language a computer can mechanically
      check (Lean, Coq, Isabelle).  The strongest available
      guarantee: subject only to the correctness of the
      definitions, the result is true.
  \end{enumerate}
\end{keyresult}

We treat the three layers in turn.

\subsubsection*{Consensus, and the event horizon of peer review}

We trust scientific papers because they are reviewed.  How
much should that trust be worth?  An order-of-magnitude
calculation done in the room is sobering.

\begin{keyresult}[The peer-review budget]\label{kr:event-horizon}
  Take the arXiv as the visible fraction of scientific
  output: $\sim40$ new submissions per day in quantum physics
  alone, $\sim10^4$ papers per year.  A genuinely careful
  referee report---the kind that would catch a non-trivial
  error---requires $3$--$4$ working days.  Honest review of
  the field's annual output thus requires $\sim3 \times 10^4$
  expert working-days per year.

  The pool of qualified experts in the relevant sub-field is
  in practice closer to $10^3$--$10^4$ than to the formal
  membership of the learned societies; sharply discipline-
  specific reviewing pushes it lower still.  We are therefore
  at, or already past, the budget at which the field can be
  honestly reviewed.  The compromise---reading the abstract,
  scanning the figures, signing off---is not a moral failing;
  it is what the budget allows.
\end{keyresult}

\begin{historical}[Voevodsky's discomfort]
  Vladimir Voevodsky was awarded the Fields Medal in 2002
  for work in motivic cohomology.  In the years following,
  he became uncertain whether his own celebrated proofs
  were correct.  The community knew of unstated subtleties
  and \emph{worked around} certain results without
  formally retracting them.  Voevodsky's response was to
  invest the rest of his career in the development of
  homotopy type theory and computer-checked mathematics.
  This is, for our purposes, the canonical example of the
  consensus layer failing on a researcher's own published
  work, and of formalisation arising as the cure.
\end{historical}

If the consensus layer is straining in 2026, layers of
verification stronger than consensus are not optional.

\subsubsection*{Computational test}

The neural-network community arrived at a parallel
realisation a decade earlier and now imposes, on its top
conferences, the requirement that submissions be accompanied
by a runnable code repository.  ``Does it compile, do the
tests pass'' is, in a certain sense, a higher bar than what a
human referee provides on a $30$-page paper: the code has
been forced through a deterministic checker.

\begin{keyresult}[Programs as referees]\label{kr:programs}
  For any conjectural result reducible to a computation
  ($N$-body simulation, special-case verification of an
  identity, brute-force search), the cheapest meaningful
  verification is to ask a coding agent to produce the
  programme, run it, and report the answer.  This is the
  mode of verification that is most accelerated by the
  technology of \S\ref{sec:agent-loop}, and it is now
  available, in practice, to anyone.
\end{keyresult}

\subsubsection*{Formalisation}

The strongest layer is mechanical proof-checking.  Lean,
together with its mathematical library Mathlib, has reached
the point where a substantial fraction of an undergraduate
mathematics curriculum---through and including parts of
functional analysis, $C^*$-algebras, category theory, and
differential geometry---is formalised, and can be invoked
rather than re-proved.  Mathlib is unique among the proof
libraries in actively chasing breadth: a small number of
people (Buzzard, Tao, Kontorovich, others) have made it
their explicit project to embed the working corpus of
mathematics into a computer-checkable form.

\begin{keyresult}[Why formalisation is now relevant to us]\label{kr:formalisation-now}
  Coding agents have been trained on Lean, Mathlib, and the
  primary literature documenting both.  They emit Lean
  proofs at a rate and with a fluency that, even very
  recently, no one anticipated.  The combination of
  (i) a strong external verifier (the Lean kernel),
  (ii) a large pre-existing library (Mathlib), and
  (iii) a fluent generator (the agent),
  closes the loop: the agent proposes, the kernel disposes,
  and the user is left only with the work of choosing
  problems and reading definitions.
\end{keyresult}

\subsection{Live demonstration: Lean with a coding agent}\label{sec:lean-demo}

The remainder of the lecture was a live demonstration in
this third regime.  The procedure was deliberately
unstructured: open a temporary directory containing a
recent Mathlib import, hand a coding agent (Claude, in this
case) a mathematical statement, and watch.

\paragraph{Warm-up: $\sqrt{2}$ irrational.}
This is the ``hello world'' of Lean.  The agent produced a
short, compiling proof in one pass; the result it actually
proved was that $2$ is not a perfect square in~$\N$, which
is the form Lean prefers in order to avoid clearing
denominators on a real square root.

\paragraph{Substantive case: tensor commutativity of line bundles.}
We then asked for the proof of a result drawn from a current
homework sheet:
\[
  L_1 \otimes L_2 \;\cong\; L_2 \otimes L_1
\]
as line bundles on a smooth manifold~$M$.  The result is
not in Mathlib in this exact form; the agent's task is to
find the right level of abstraction (Mathlib's category-
theoretic infrastructure for vector bundles being
substantial) and to emit a proof that compiles.

The behaviour observed was characteristic.  The agent first
located \texttt{Mathlib.Topology.VectorBundle} and similar
files, complained about the size of the project (``500 to
$1{,}500$ lines, four pieces, all mirroring existing
patterns''), wrote a proof skeleton with the missing pieces
filled by the placeholder \texttt{sorry}, and began an
iterative loop of attempting each placeholder, observing
the type-check error, and adjusting.  Within a session the
proof compiled; the resulting code was, by the agent's own
description, \emph{packaging rather than new mathematics}.

\begin{intuition}[The role of \texttt{sorry}]
  The \texttt{sorry} placeholder is what makes large-scale
  agent-driven formalisation tractable.  It allows a proof
  skeleton to type-check before any of its leaves have been
  proved; the agent can then descend the tree, replacing
  one \texttt{sorry} at a time, and each individual proof
  obligation is small enough to be handled.  Decomposition,
  in this domain, is not a metaphor: it is a primitive of
  the language.
\end{intuition}

\begin{keyresult}[An empirical data point]\label{kr:80-pages}
  An honest figure of present capability: with a current
  frontier model, a domain expert can, in a weekend,
  produce of order $80$ pages of \LaTeX\ together with
  formalised Lean proofs in a sub-area such as functional
  analysis or operator algebras.  The bottleneck is no
  longer the formalisation; it is the rate at which the
  human can supply ideas.
\end{keyresult}

\subsection{Failure modes of formalisation}\label{sec:formalisation-failures}

Several failure modes deserve to be named explicitly.

\paragraph{Wrong definitions.}
Lean checks proofs against the definitions one writes, not
against the definitions one \emph{intends}.  An off-by-an-axiom
definition (\emph{``a smooth manifold is a manifold where
the transition functions are twice differentiable''})
admits proofs that are mechanically valid but
mathematically vacuous.  Definitions must be read.  In
practice the worst cases self-disclose: a faulty definition
quickly produces consequences that are visibly absurd, and
the contradiction surfaces before the formalisation grows
large.  This is a useful but unprovable heuristic.

\paragraph{Wheel-spinning on false statements.}
When asked to prove a statement that is in fact false, the
agent does not, in general, find the counter-example.  It
recurses, fails, re-tries, and consumes turns.  Counter-
example search is a separate skill; for that, a Python
program is usually faster.  The takeaway is workflow:
\emph{always run the counter-example pass before launching
the formalisation pass}.

\paragraph{Context-window saturation.}
The proof of a single non-trivial result---in the live
demo, the line-bundle commutativity---consumed of order
$10^5$ tokens of working context.  Even a model with a
$10^6$-token context window cannot, in practice, hold more
than a handful of such proofs simultaneously while
maintaining inference quality.  A long programme of
formalisation must therefore be \emph{decomposed} so that
each subgoal fits independently in context.  This is the
same lesson as the $P$ vs $\textsf{NP}$ discussion of
\S\ref{sec:ambitious-science}, restated for proof
assistants.

\paragraph{Erd\H{o}s slop.}
A failure mode at the level of community, not of an
individual run.  Erd\H{o}s posed thousands of conjectures,
most of which have never been seriously attacked because
the supply of conjectures vastly exceeded the supply of
mathematicians interested in them.  It is presently
fashionable to point a chat model at the Erd\H{o}s list and
publish what falls out.  Some of these claims will be
correct; verifying them, however, falls back on the same
small pool of human experts whose attention is already
oversubscribed.  This converts a free resource (compute)
into an expensive one (expert-hours), in the wrong
direction.

\begin{remark}[On Erd\H{o}s slop]
  The structural problem is asymmetric: an LLM that emits
  ten thousand candidate Lean proofs of Erd\H{o}s problems
  taxes the verification budget of exactly the experts who
  could otherwise be doing science.  ``Slop'' is the
  appropriate word; the polite version is that we are
  generating positive externalities for ourselves while
  imposing negative ones on the human community of record.
\end{remark}

\subsection{The singularity, revisited}\label{sec:singularity}

The discussion turned, by way of an audience question, to
the recursive-self-improvement thesis: the claim that an AI
capable of editing its own substrate will iteratively
amplify its own intelligence to a runaway, qualitative
threshold (``AGI,'' ``the singularity'').  The view of these
seminars (\S\ref{sec:self-improvement}) is that this thesis
conflates two different objects.  The harness, written in
plain code, is rewritten at every turn; the weights, in
which the cognition lives, are not.  The cognitive ceiling
is set by the weights and is therefore stationary on
inference-time scales.

A sharper formulation of the objection emerged.

\begin{keyresult}[Recursive self-improvement as an optimisation problem]\label{kr:singularity}
  The singularity claim, made concrete, requires a process
  that climbs from one local optimum of the
  ``intelligence-as-a-function-of-parameters'' landscape to
  a strictly better one.  But this is exactly the global-
  optimisation problem that all of optimisation theory tells
  us is hard.  The state required to even \emph{compare}
  candidate optima---a record of which configurations have
  been tested---grows with the dimension of the parameter
  space, and one can write down dimension counts at which
  the comparison memory exceeds, by many orders of magnitude,
  any conceivable physical store.  In short: \emph{we run
  out of memory before we run out of local minima}.
\end{keyresult}

\begin{remark}[Status of the argument]
  This sketch is in the spirit of a thermodynamic
  no-go theorem rather than a fully rigorous result.  The
  components---bumpy landscapes, the inadequacy of greedy
  ascent, dimensional explosion of search space---are
  individually well-understood.  Assembled into a single
  ten-step deductive argument they would constitute, in my
  view, a genuine obstruction to the strong singularity
  thesis.  Whether and how to write that argument down is
  one of the candidate problems for this seminar.
\end{remark}

\begin{audience}[Audience reaction]
  The question that prompted the discussion was the
  honest one: \emph{``I would like to be convinced that AGI
  is not a threat.''}  The honest reply: I cannot prove
  this, because the term is not well-defined enough to be
  the object of a proof; but the recursive-self-improvement
  argument, which is the strong technical version of the
  threat narrative, has the structure of a global-
  optimisation claim and inherits the well-known
  difficulties of such claims.  This is not a guarantee.
  It is a reason to be sceptical of the strongest version
  of the threat.
\end{audience}

\subsection{Breadth, depth, and the visualisation imperative}\label{sec:breadth-visualisation}

The final third of the seminar returned to the operational-
competency thesis of \S\ref{sec:ambitious-science}, with a
sharper claim.

\begin{keyresult}[Breadth as the new affordance]\label{kr:breadth}
  Becoming a domain expert---a process that inevitably
  narrows one's competence outside the domain---used to be
  the only path to working knowledge of a field.  With
  effective coding agents and good verification, a worker
  with operational competency
  (Definition~\ref{def:opcomp}) can reach \emph{useful}
  competence in many neighbouring domains in days rather
  than years.  The corresponding research mode is one of
  \emph{breadth}: the previously prohibitive cost of
  entering an adjacent field has fallen sharply.
\end{keyresult}

The limitation, naturally, is depth.  Deep insight remains
slow.  But for many of the questions a physicist actually
wants to ask---``does this idea make contact with that
field?'', ``is this construction equivalent to that one?'',
``what does this look like in three dimensions?''---breadth
is what was missing, and breadth is what we now have.

A particular instance of this affordance was developed at
length: \emph{visualisation as an instrument of
understanding}.

\begin{intuition}[Visualisation as a question-answering channel]
  Sitting in a lecture, one's understanding accumulates a
  small collection of \emph{spinning plates}: questions
  that arise but cannot be asked without disrupting the
  pace of delivery, and that consume cognitive cycles for
  the rest of the lecture until one of them resolves.  A
  personal assistant that answers each question instantly,
  and where appropriate by producing an animation or an
  interactive demo, removes the spinning plates and frees
  the cycles.  The pedagogy of a textbook---linear,
  non-interactive, one-pass---is precisely what blocks
  this.
\end{intuition}

\begin{historical}[A worked example: hydrogen, helium, neutron stars]
  In the spirit of demonstrating the claim, three
  visualisations were prepared with help from a coding agent.
  \begin{itemize}
    \item A hydrogen atom in a~$2p$~state, in an external
      magnetic field, undergoing spontaneous emission to
      the~$1s$ ground state.  The plotted quantity is the
      true electron density (computed from the density
      matrix, not a pure-state stand-in).  This is the
      picture I had wanted on the blackboard for a decade.
    \item Helium scattering, with the electron probability
      cloud bouncing between the nuclei.
    \item The optical spectrum of our nearest neutron star,
      with the relativistic redshift and the squeezed
      hydrogenic wave-functions appropriate to the
      colossal magnetic field of the atmosphere.
  \end{itemize}
  Each was the work of an afternoon.  None would have
  been undertaken without a coding agent at hand; each, on
  reflection, accelerates one's intuition more than a
  semester of equations.
\end{historical}

\begin{remark}[On the asymmetry of teaching and learning]
  A second-order observation: the affordance is not
  symmetric between teachers and learners.  A teacher
  builds visualisations to communicate \emph{depth} in a
  domain they already understand; a learner uses them to
  acquire \emph{breadth} across domains they do not yet.
  Both are valuable, and the technology serves both, but
  the workflows are different.  A learner who feels
  pressure to acquire depth before breadth has the
  ordering wrong for the present moment.
\end{remark}

\subsection{The mini-project: a question for next week}\label{sec:mini-project}

The seminar closed by re-stating, with somewhat more force,
the request of \S\ref{sec:ambitious-science}: each
participant should bring, by the next session, a concrete
formulation of \emph{the world they would like to create}.

\begin{keyresult}[The non-technical homework]\label{kr:homework-world}
  The technical bottleneck has moved.  The new bottleneck is
  the choice of problem.  The exercise---which is harder
  than it sounds---is to articulate, in a sentence or two, a
  research-scale question one would actually like to answer,
  on the assumption that operational competency in any
  required adjacent field is now cheap.  We will work
  outwards from these statements in the remaining sessions.
\end{keyresult}

\begin{audience}[Examples raised in the room]
  Several candidate directions surfaced:
  \begin{itemize}
    \item the energy and monetary cost of frontier-model
      inference, as a quantitative survey
      (cf.\ the third homework of \S\ref{sec:bitter-lesson});
    \item a precise statement of Bell's theorem and a
      reconciliation of the various extant formulations,
      in the spirit of ``which version is true?'';
    \item three-dimensional, physically-correct
      visualisations of canonical quantum processes
      (\S\ref{sec:breadth-visualisation});
    \item more speculatively, a written-out form of the
      anti-singularity argument of
      Theorem~\ref{kr:singularity}.
  \end{itemize}
  None of these is final.  All are, on present terms,
  attemptable.
\end{audience}

\begin{remark}[The shape of the remaining series]
  Whether the seminar develops as a sequence of independent
  mini-projects, a single coordinated attack on a larger
  question, or a structured exploration of a small number of
  candidate worlds, is for the participants to decide.  The
  prepared notes (\S\ref{sec:agent-loop}--\S\ref{sec:ambitious-science})
  exist as a resource; the live work begins from here.
\end{remark}
